\clearpage
\appendix

\section{Experimental Setup Details}
\label{app:setup}

\subsection{Task and Dataset}
\label{app:dataset}
The HuggingFace \texttt{stanfordnlp/sst2} train set (67,349) consists of phrases, so we split at the sentence level(Table~\ref{tab:app-splits}).

We use the train set for fitting probes, SHAP/IG/GA background and baseline, and steering $\alpha$ calibration. The val set is used for probe hyperparameters, the activation-frequency feature mask, and reported probe accuracy. The held-out split is used for all attribution (SHAP/IG/GA) and all causal evaluation (ablation/steering).

\begin{table}[h]
\centering
\caption{Data split sizes and purposes.}
\label{tab:app-splits}
\begin{tabular}{l r}
\toprule
Split & $n$ \\
\midrule
Probe train & 6{,}696 \\
Probe val & 1{,}435 \\
Held-out (``shap'') & 1{,}436 \\
\bottomrule
\end{tabular}
\end{table}

\paragraph{Canonical row sampling.} All evaluation and background rows are
\emph{prefixes} of one seeded permutation (seed$=42$), not independent
draws. The 100-row local sample is therefore exactly the first 100 rows of
the 1000-row global sample, so local and global results join by row
position, and the 100-row train background does not depend on evaluation
size. Individual experiments may use a prefix of the 1000 rows (100 for
local per-example runs).

\subsection{Model and Probes}
\label{app:model}
\label{app:probes}

GPT-2 Small and the layer-7 \texttt{resid\_post} SAE are frozen throughout.
SAE decoder rows are approximately unit norm (median 0.9996 across all
32{,}768 features), so we apply no decoder-norm correction to any steering
vector.

\paragraph{Linear probe.} \texttt{sklearn} \texttt{LogisticRegression}
with an L1 penalty ($C=1$, \texttt{liblinear} solver, fixed
\texttt{random\_state}) over the full 32{,}768-d activation. The number of
non-zero coefficients is run-dependent ($\sim$600), so the 114-feature
mask derived from it is a property of the reported run, not a fixed
constant.

\paragraph{MLP probe.} \texttt{sklearn} \texttt{MLPClassifier} with one
hidden layer of width 32, ReLU activation, Adam, and L2 regularization,
trained on the 114 masked features only.

\clearpage
\section{Attribution Method Details}
\label{app:attribution}

\paragraph{Global SHAP.} We report
$\Phi_j = \mathbb{E}_e[\phi_j(x_e)]$, the \emph{mean signed} SHAP value per
feature over the 1000-row held-out sample. Sign is preserved, and $\Phi$ is
the basis for every directional comparison.

\paragraph{Integrated Gradients.}
\label{app:ig}
\[
\mathrm{IG}_j = (x_j - b_j) \cdot \int_0^1 \frac{\partial f}{\partial x_j}\big(b + \alpha(x-b)\big)\, d\alpha,
\]
approximated in the MLP arm with a 50-step trapezoidal Riemann sum. IG and
GA share the baseline $b$, the \emph{mean} of the 100-row train background
rather than a zero vector. SAE features very sparse, so a zero
baseline would make GA degenerate to 0 whenever $x_j=0$. For the linear
probe, IG, GA, and interventional SHAP all reduce to
$w_j \cdot \mathbb{E}_e[x_{e,j} - \mu_j]$; we verify this numerically
(Appendix~\ref{tab:app-identity}).

\clearpage
\section{Causal Evaluation Details}
\label{app:causal}

\subsection{Steering-Strength ($\alpha$) Calibration}
\label{app:alpha}
We sweep $\{0.5,1,2,4\}\times$ the median natural positive-activation scale
over three top-$|\mathrm{SHAP}|$ pilot features on 200 train examples,
reading out a layer-11 residual probe's $P(\text{positive})$ rather than
the evaluation logit difference. $\alpha$ is meant as a rough
``measurable, unsaturated'' magnitude, not a hyperparameter tuned against
the evaluation metric. We select the smallest $\alpha$ with
$\text{mean}|\Delta P| \geq 0.02$ and newly saturated fraction (examples
with $P<0.02$ or $P>0.98$ after steering that were not saturated before
steering) $\leq 0.25$, giving $\alpha=0.4936$
(Table~\ref{tab:app-alpha-pilot-raw}).

\begin{table}[h]
\centering
\caption{Raw $\alpha$-calibration pilot sweep (probe-readout, 3 pilot features, 200 train examples).}
\label{tab:app-alpha-pilot-raw}
\begin{tabular}{cccc}
\toprule
$\alpha$ & mean$|\Delta P|$ & mean signed $\Delta P$ & newly saturated frac.\ \\
\midrule
0.4936 & 0.0490 & $-0.0301$ & 0.028 \\
0.9871 & 0.0961 & $-0.0571$ & 0.085 \\
1.9743 & 0.1791 & $-0.0981$ & 0.182 \\
3.9485 & 0.2969 & $-0.1422$ & 0.368 \\
\bottomrule
\end{tabular}
\end{table}

\paragraph{Per-feature scaling.} In the primary \textbf{scaled} mode,\\
$\alpha_i = \alpha \cdot \mathrm{scale}_i / \mathrm{median}(\mathrm{scale})$,
where $\mathrm{scale}_i$ is feature $i$'s mean positive train activation.
This equalizes the relative nudge across features whose natural scales
differ severalfold. The \textbf{constant} mode (flat $\alpha$, robustness
check) confounds the $|\text{attribution}|$-vs-$|\Delta|$ correlation,
since $|\Delta|$ then partly reflects how hard a feature was pushed
relative to its own scale.

\subsection{Candidate Selection}
\label{app:selection}

\texttt{top} selection takes the top-$k$ by $|\mathrm{SHAP}|$ and scores all
methods on that set. This range-restricts only SHAP's scored variance,
which mechanically deflates SHAP's $\rho$ relative to the other methods.
\texttt{all} scores every masked feature ($n=114$). At $n=20$ the
confidence interval on Spearman $\rho$ spans most of $[-1,1]$, so a single
$k=20$ random draw cannot distinguish a null from a moderate effect.

\paragraph{Provenance of the \texttt{all} analysis.} The \texttt{all}
selection was introduced after an earlier analysis of this work (the
\emph{initial analysis}, on a phrase-level split we later found leaked
sub-phrases across splits), whose single $k=20$ draw looked null for several
tests. An earlier version of this manuscript presented \texttt{all} as the
powered result after seeing that $k=20$ outcome, which is outcome-dependent.
We have since rebuilt the split at sentence level and regenerated every
number reported here (the \emph{final analysis}). The \texttt{all}
selection, the BH family, and the rule for it were fixed before the final
analysis was run, so it is a confirmatory test of a design chosen on the
initial data, not a second look at the same data. All $k=20$ results from the final analysis are
reported (Table~\ref{tab:headline-faithfulness}, and the multi-seed sweep below). We treat
\texttt{all} as the inferential run because it uses the whole filtered pool
with no sampling, not because of how $k=20$ came out: in the final analysis the $k=20$
draw is also positive for importance and null for direction, so the two
agree.

\subsection{Multiple-Comparison Correction}
\label{app:bh}

\paragraph{Family rule.} We apply Benjamini--Hochberg (BH, FDR 0.05) with the
following primary family: the tests of one \texttt{all} run (one arm, one
intervention scope; 8 linear, 7 MLP), and, for local faithfulness, the four
crossed-bootstrap tests (Block A/B $\times$ importance/directional) of one
local run. The rule is a property of how the experiments are run, not of
their outcomes. The per-run correction was specified before the final
analysis was run, but after the initial analysis had been examined. The wider
families below were added after the final results were in and are post hoc. We report them because the choice of family is the
author's, so we show whether any verdict depends on it.

\paragraph{Duplicate tests.} In the linear arm SHAP, IG, and GA have identical
$(\rho, p)$ by construction (the three methods are analytically the same for a
linear probe), so the 8 linear tests are 4 distinct ones. Counting them
separately inflates survivors and, when pooled, BH power. We therefore report
the main sensitivity analysis with identical tests counted once (22 distinct
global tests); the raw counts are 16/30.

\paragraph{Sensitivity to the family.} Counts are recomputed from the per-test $p$-values of every inferential test in the paper under each family definition.
\begin{table}[h]
\centering
\small
\caption{Tests surviving FDR 0.05 under alternative family definitions. BY is
Benjamini--Yekutieli (arbitrary dependence).}
\begin{tabular}{lcc}
\toprule
Family & BH & BY \\
\midrule
Global, per run (primary) & 12/22 & 11/22 \\
Global, pooled over arms, scopes, or both & 12/22 & 11/22 \\
Global, raw (IG/GA counted separately) & 16/30 & 15/30 \\
\midrule
Paper-wide, per run / per local source (primary) & 13/32 & 11/32 \\
Paper-wide, global and local pooled separately & 12/32 & 11/32 \\
Paper-wide, one family of all inferential tests & 12/32 & 11/32 \\
\bottomrule
\end{tabular}
\end{table}
The paper-wide family has 32 tests: the 22 distinct global tests and the 10
local crossed-bootstrap tests (three local runs; the language-model
diagnostic has importance blocks only). The sign tests, minimal-pair
summaries, and seed sweep carry no inferential claim in the paper and are
excluded.

Two tests have a verdict that depends on the family.
\begin{itemize}
\item Last-token MLP SHAP importance survives under BH in every family
($q=0.030$ in its own run, $0.034$ pooled) but not under BY ($q=0.14$).
\item The language-model-target local diagnostic, Block A importance, survives
within its own local family ($q=0.048$) but not when pooled with the other
tests ($q=0.055$).
\end{itemize}
Every other verdict is the same under every family. Importance faithfulness of
SHAP, IG, and GA (all-token, both arms; last-token linear and IG/GA in MLP)
survives all families including BY, and no directional test of SHAP, IG, or GA
has $q<0.25$ in any. No linear- or MLP-arm local test survives in any family. A
test that fails the FDR criterion is not evidence that its effect is zero.

\paragraph{Dependence.} The two scopes share features and eval rows, and the
local runs share sentences, so the tests are not independent. BH is valid under
positive dependence, which we think is plausible here, but we do not
verify it; the BY column does not assume it.

\paragraph{Seed sweep.} The multi-seed sweep uses a different convention. Its
10 seeds $\times$ 8 (linear) or 7 (MLP) per-seed tests each use $n=20$
features that overlap across seeds, so the pooled BH count over them (80
linear tests, 70 MLP) is descriptive only.

\subsection{Local Faithfulness}
\label{app:local}

Local analyses use $k_{\text{local}}=k_{\text{global}}=20$ and
$n_{\text{examples}}=100$ in both arms. We report Block B (active features
only), not Block A, as the headline local-vs-global comparison, since
Block A is dominated by the tied block of inactive features.

\clearpage
\section{Multi-Feature (Group) Steering Details}
\label{app:group}

\begin{itemize}
  \item \textbf{Vector construction}: for method $m$, we take the top-$k$
  features by $|\text{score}|$ from the shared 114-feature pool and set
  $v = \sum_i \text{score}_i \cdot W_{\text{dec}}[i]$, unit-normalized.
  \item \textbf{\texttt{actdiff} baseline}: top-$k$ by
  $|\overline{a}^{+} - \overline{a}^{-}|$ (mean positive minus mean
  negative activation), weighted by the signed difference. It uses no
  probe, so it is scored once and shown in both arms.
  \item \textbf{No MLP probe curve}: MLP probe saliency is unsigned, and
  summing unsigned weights would add positive- and negative-sentiment
  directions with the same sign. The logistic arm's probe curve is kept
  because its coefficients are signed.
  \item \textbf{Effect and cost}: effect is mean $\Delta$ logit difference
  (wonderful $-$ awful); cost is mean $\Delta$ perplexity against the clean
  baseline.
  \item \textbf{Baseline}: we add $\alpha \cdot v$ directly to the clean
  residual stream with no SAE encode/decode, so the baseline is a clean
  forward pass rather than the per-feature experiments' reconstruction.
  Effect sizes are therefore not comparable across the two families; we
  compare only curve shapes.
\end{itemize}

\section{Minimal-Pair Stimuli}
\label{app:minimal-pairs}

\texttt{remove} pairs take held-out sentences containing \emph{exactly one}
word from a hand-curated sentiment lexicon and delete that token, leaving
everything else byte-identical. With two sentiment words, the remaining one
would keep carrying sentiment and the edit would no longer have a single
cause. \texttt{insert} pairs append ``\,it was
\textlangle word\textrangle\,.'' to lexicon-free sentences.
\texttt{substitute} pairs swap that word in place for a random word of
opposite polarity. \texttt{insert-interior} pairs place ``it was
\textlangle word\textrangle,'' at the midpoint of a lexicon-free sentence,
so the final token and its immediate context are unchanged. Only 5\% of
\texttt{insert} pairs change the identity of the final token (both
sentences end in ``.''), so the type/position confound concerns distance
from the final token, not a replaced token.

\paragraph{Changed-feature threshold.} A candidate counts as changed if
$|\Delta a_j|\ge 0.10\max_i|\Delta a_i|$ within the pair. Local-SHAP AUROC
(random distractors) is stable across relative thresholds $0.05$--$0.50$ for
insertion ($0.62$--$0.68$) but not for removal ($0.55$ at $0.10$, $0.40$ at
$0.50$); absolute thresholds $0.1$--$1.0$ give $0.61$--$0.62$ for insertion and a non-monotone $0.38$--$0.56$ for removal. A
threshold-free Spearman of $|$score$|$ against $|\Delta a|$ per pair gives
local SHAP $0.04$ (removal) and $0.26$ (insertion), against $0.57$ and $0.67$
for activation magnitude at the edited last token.

\paragraph{Readout noise floor.} Null edits (a neutral word deleted or
appended) give a median $|\Delta|$ of $0.14$ logits (removal-type) and $0.23$
(insertion-type), against mean effects of $0.28$ and $0.77$. Restricting to
pairs above the null 90th percentile leaves 26 removal pairs (sign accuracy
$0.85$) and 52 insertion pairs ($0.94$); AUROC and $\rho$ on those subsets are
within $0.03$ and $0.06$ of the full-sample values, with the same method ordering.

\paragraph{Seeds.} Seeds $0$--$3$ ($0$--$2$ for \texttt{substitute}) redraw pairs
and distractors from the same held-out pool, so they overlap and are descriptive.
Per-edit standard deviations of local-SHAP AUROC under random distractors are
$\le 0.03$, and no local or global SHAP run exceeds $0.5$ under activity-matched
distractors.

\clearpage
\section{Supplementary Tables and Figures}
\label{app:supptables}

\subsection{Linear-Arm Identity Check and MLP-Arm Exact $p$-Values}

\begin{quote}
\textit{Closed-form identity check (linear probe, margin target):}
$\max|\mathrm{IG}-\Phi| = 1.72\times10^{-7}$ (relative $1.75\times10^{-6}$),
$\max|\mathrm{GA}-\Phi| = 1.72\times10^{-7}$. IG, GA, and
interventional SHAP coincide analytically for a linear probe; we obtain
pairwise $\rho = 1.000$ ($n=615$ non-zero-coefficient features on this
fit) as an identity, not an empirical result.
\end{quote}

\begin{table}[h]
\centering
\caption{MLP-arm attribution agreement (Spearman $\rho$ on $|\text{score}|$, $n=114$).}
\label{tab:app-mlp-ranking}
\begin{tabular}{lcc}
\toprule
Pair & $\rho$ & $p$ \\
\midrule
Probe vs.\ IG   & 0.287 & 0.002 \\
Probe vs.\ SHAP & 0.250 & 0.007 \\
Probe vs.\ GA   & 0.242 & 0.010 \\
IG vs.\ SHAP    & 0.866 & $<0.0001$ \\
IG vs.\ GA      & 0.862 & $<0.0001$ \\
SHAP vs.\ GA    & 0.693 & $<0.0001$ \\
\bottomrule
\end{tabular}
\end{table}
\label{tab:app-identity}

\subsection{Selection-Bias and Sign-Stratified Checks}
\label{app:selection-table}

\begin{table}[h]
\centering
\caption{SHAP faithfulness across candidate-selection strategies (steering), $k=20$ except \texttt{all} ($k=114$). \texttt{top} biases SHAP's $\rho$ downward relative to other methods; \texttt{top-pos}/\texttt{top-neg} test the sign-inversion hypothesis without mixed-sign cancellation.}
\label{tab:app-selection-variants}
\resizebox{\columnwidth}{!}{%
\begin{tabular}{lccccc}
\toprule
Arm : kind & random & top & top-pos & top-neg & all ($k$=114) \\
\midrule
Linear : imp & $+0.555^{*}$ & $+0.496^{*}$ & $+0.293$ & $+0.397$ & $+0.377^{**}$ \\
Linear : dir & $-0.077$ & $-0.083$ & $+0.268$ & $-0.226$ & $-0.097$ \\
MLP : imp    & $+0.744^{**}$ & $+0.107$ & $+0.322$ & $+0.648^{**}$ & $+0.347^{**}$ \\
MLP : dir    & $-0.147$ & $-0.161$ & $+0.168$ & $-0.376$ & $-0.091$ \\
\bottomrule
\end{tabular}%
}
\vspace{2pt}

{\footnotesize $^{*}p<0.05$, $^{**}p<0.01$, uncorrected; all other cells $p>0.08$.}
\end{table}

\subsection{Alpha-Mode Robustness (Scaled vs.\ Constant)}
\label{app:alpha-mode-table}

\begin{table}[h]
\centering
\caption{Alpha-mode sensitivity: scaled vs.\ constant per-feature $\alpha$ (steering, random selection, $k=20$).}
\label{tab:app-alpha-mode}
\resizebox{\columnwidth}{!}{%
\begin{tabular}{llcc}
\toprule
Arm & Method & $\rho_{\text{imp}}$ (scaled / constant) & $\rho_{\text{dir}}$ (scaled / constant) \\
\midrule
Linear & SHAP/IG/GA & $+0.555$ / $+0.194$ & $-0.077$ / $+0.072$ \\
Linear & Probe      & $+0.519$ / $+0.226$ & $+0.624$ / $+0.567$ \\
MLP    & SHAP       & $+0.744$ / $+0.459$ & $-0.147$ / $-0.011$ \\
MLP    & Probe saliency & $+0.400$ / $+0.365$ & n/a \\
\bottomrule
\end{tabular}%
}
\end{table}

A large gap here would indicate that constant-$\alpha$ steering numbers
substantially reflect activation scale rather than faithfulness; we report
both modes regardless of which way the comparison comes out.

\subsection{Multi-Seed Replication}
\label{app:seed-sweep}

The sweep is descriptive. Seeds draw overlapping subsets of one
114-feature pool and reuse the same per-feature effects, so they are not
independent replicates; inference comes from the \texttt{all} run
(Appendix~\ref{app:selection}). Few per-seed tests survive a pooled BH
correction (0/80 linear, 3/70 MLP; the three MLP survivors are seed 0's
SHAP, IG, and GA importance tests). This is expected because each per-seed
test uses only $n=20$ features. In the $n=114$ \texttt{all} run, the same
importance effects survive under every BH family definition
(Appendix~\ref{app:bh}). The sweep's mean importance $\rho$ is positive in
every method--arm cell, and SHAP/IG/GA directional $\rho$ is centered near
zero, matching the \texttt{all} run.

\begin{table}[h]
\centering
\small
\setlength{\tabcolsep}{4pt}
\caption{Multi-seed faithfulness sweep (10 seeds, $k=20$, scaled $\alpha$). Descriptive only. Ref $p$ is the fraction of seeds with $|\rho|$ at least that of the published seed-0 draw.}
\label{tab:app-seed-sweep}
\begin{threeparttable}
\begin{tabular}{llrrrrrr}
\toprule
Method & Type & Mean $\rho$ & SD & Min & Max & Frac.\ $>0$ & Ref $p$ \\
\midrule
\multicolumn{8}{l}{\textit{Linear arm}} \\
SHAP/IG/GA & imp & $+0.339$ & 0.199 & $-0.032$ & $+0.632$ & 0.90 & 0.200 \\
SHAP/IG/GA & dir & $-0.047$ & 0.160 & $-0.356$ & $+0.141$ & 0.50 & 0.700 \\
Probe      & imp & $+0.236$ & 0.268 & $-0.260$ & $+0.582$ & 0.80 & 0.200 \\
Probe      & dir & $+0.422$ & 0.162 & $+0.162$ & $+0.624$ & 1.00 & 0.100 \\
\midrule
\multicolumn{8}{l}{\textit{MLP arm}} \\
SHAP       & imp & $+0.368$ & 0.189 & $+0.128$ & $+0.744$ & 1.00 & 0.100 \\
SHAP       & dir & $-0.005$ & 0.191 & $-0.377$ & $+0.202$ & 0.50 & 0.700 \\
GA         & imp & $+0.460$ & 0.125 & $+0.316$ & $+0.752$ & 1.00 & 0.100 \\
\bottomrule
\end{tabular}
\begin{tablenotes}\footnotesize
    \item Imp: importance; dir: directional. Pooled BH (descriptive): 0/80 tests survive (linear), 3/70 (MLP), at FDR 0.05.
\end{tablenotes}
\end{threeparttable}
\end{table}




\subsection{Single-draw $k=20$ causal validation results}
\label{app:k-20}

\begin{table}[H]
\centering
\small
\setlength{\tabcolsep}{5pt}
\caption{Single-draw causal faithfulness ($k=20$, steering and ablation, random selection, $n_\text{eval}=1000$). Cells show Spearman $\rho$ with uncorrected $p$ in parentheses; at $n=20$ these are descriptive and not used for inference. For the linear arm, IG and GA are analytically identical to SHAP and omitted as separate rows.}
\label{tab:headline-faithfulness}
\begin{tabular}{lcc}
\toprule
Method & $\rho_\text{imp}$ ($p$) & $\rho_\text{dir}$ ($p$) \\
\midrule
\multicolumn{3}{l}{\textit{Linear arm, steering}} \\
SHAP/IG/GA     & $+0.555$ (0.011) & $-0.077$ (0.748) \\
Probe          & $+0.519$ (0.019) & $+0.624$ (0.003) \\
\addlinespace
\multicolumn{3}{l}{\textit{Linear arm, ablation}} \\
SHAP/IG/GA     & $+0.511$ (0.021) & $-0.427$ (0.060) \\
Probe          & $+0.347$ (0.133) & $+0.648$ (0.002) \\
\midrule
\multicolumn{3}{l}{\textit{MLP arm, steering}} \\
SHAP           & $+0.744$ (0.0002) & $-0.147$ (0.535) \\
Probe saliency & $+0.400$ (0.081) & n/a \\
IG             & $+0.749$ (0.0001) & $-0.146$ (0.540) \\
GA             & $+0.752$ (0.0001) & $-0.171$ (0.470) \\
\addlinespace
\multicolumn{3}{l}{\textit{MLP arm, ablation}} \\
SHAP           & $+0.568$ (0.009) & $-0.408$ (0.075) \\
Probe saliency & $+0.281$ (0.230) & n/a \\
IG             & $+0.573$ (0.008) & $-0.421$ (0.065) \\
GA             & $+0.600$ (0.005) & $-0.453$ (0.045) \\
\bottomrule
\end{tabular}
\end{table}


\paragraph{Continuous (partial-Spearman) activity control.} We compute a
partial Spearman of $|\text{attribution}|$ vs.\ $|\Delta|$ controlling
continuously for $|\text{activation}|$ (rather than the binary
active/inactive split of Block B), keeping the whole candidate set. Local
SHAP's mean partial $\rho$ is near zero ($+0.038$ linear / $+0.063$ MLP),
versus $-0.168$ / $-0.191$ for global $\Phi$.


\subsection{Group-Steering Full Sweep}
\label{app:group-full}

This is the full parameter grid underlying the $k=20$
spot-check: $k \in \{5,10,20,30,50,65\} \times \alpha \in
\{0.5,1,2,5,10,20,50\}$ for 7 vectors (4 logistic, 3 MLP) plus
\texttt{actdiff}.

\begin{figure}[t]
  \centering
  \begin{subfigure}[b]{\linewidth}
    \includegraphics[width=\linewidth]{figures/effect_vs_cost_logistic_grid.png}
    \caption{Logistic arm.}
  \end{subfigure}
  
  \vspace{1em} % Adds vertical space between the stacked subfigures
  
  \begin{subfigure}[b]{\linewidth}
    \includegraphics[width=\linewidth]{figures/effect_vs_cost_mlp_grid.png}
    \caption{MLP arm.}
  \end{subfigure}
  \caption{Full group-steering effect-vs-cost grid ($k\times\alpha$), one panel per method per arm. This experiment's baseline is a clean forward pass (no SAE encode/decode), unlike the per-feature steering/ablation experiments' encode$\to$decode-reconstruction baseline (Appendix~\ref{app:group}); we compare curve \emph{shapes} across arms, never magnitudes.}
  \label{fig:app-group-steering-grid}
\end{figure}

\subsection{Lexical Simplicity Control}
\label{app:lexical}

We use the full 1{,}436-row held-out split as the population here, not
the 1000-row canonical eval sample used everywhere else in the paper.

\begin{table}[h]
\centering
\caption{Lexical simplicity control (held-out, $n=1{,}436$). Models are fit on train with regularisation chosen on validation. Diff.\ is probe minus baseline accuracy with a paired-bootstrap 95\% interval; $p$ is McNemar.}
\label{tab:app-lexical-simplicity}
\begin{tabular}{lcccc}
\toprule
Model & Features & Accuracy & Diff.\ [95\% CI] & $p$ \\
\midrule
SAE probe (held-out) & 114 masked / 32768 & 0.810 & -- & -- \\
Lexicon-count LR & 158 & 0.611 & $+0.198$ [$+0.166$, $+0.229$] & $<10^{-30}$ \\
Unigram BoW LR (binary) & 13{,}605 & 0.794 & $+0.016$ [$-0.010$, $+0.041$] & 0.26 \\
Unigram+bigram BoW LR & 79{,}812 & 0.800 & $+0.010$ [$-0.015$, $+0.034$] & 0.49 \\
TF-IDF unigram+bigram LR & 79{,}812 & 0.806 & $+0.003$ [$-0.022$, $+0.028$] & 0.83 \\
Dense layer-7 residual LR & 768 & 0.870 & $-0.060$ [$-0.080$, $-0.040$] & $<10^{-8}$ \\
\bottomrule
\end{tabular}
\end{table}

The best validation-selected bag-of-words (TF-IDF) is within $0.03$ accuracy
of the probe, so by our pre-fixed rule the probe is no better than word
presence on this split. Co-activation is unchanged from the earlier analysis
(24.33 active filtered features per example; mean pairwise Jaccard of the
top-20 features $0.223$).

\paragraph{Word-detector tests.} Held-out $R^2$ of a ridge regression from
bag-of-words to feature activation: for the top-20 $|\Phi|$ features the
median is $0.171$ with full vocabulary and $0.003$ with the lexicon only; for
the other 94 masked features it is $0.106$ (90th percentile $0.287$) and
$0.003$. Four of the 20 top features reach $R^2\ge0.25$. Full bag-of-words
explains $0.389$ of the probe logit's variance ($0.148$ with the lexicon only);
this is a lower bound on reducibility, since a linear ridge fit on 6.7k
sentences with a 13.6k-word vocabulary is limited. Signed and position-aware
correlations with the lexicon are small for the top-20 features: mean $|r|$ is
$0.067$ (has-positive), $0.056$ (has-negative) and $0.049$ (lexicon word in
the last three tokens, present in 12\% of sentences); mean $|\rho|$ with the
net lexicon score is $0.076$, and the mean sign-agreement of $\Phi$ with it is
$-0.006$. The earlier lexicon-presence correlation (mean $|r|=0.048$) was
therefore a weak test: the lexicon explains almost none of any feature's
variance, so it could not have revealed word-detection.

\subsection{All-token global faithfulness}
\label{app:all-token-global-faithfulness}
\begin{table}[h!]
    \centering
    \small
    \setlength{\tabcolsep}{5pt}
    \caption{All-token global faithfulness ($k=114$, steering, all filtered features). Cells show Spearman $\rho$ with BH-corrected $q$ in parentheses.}
    \label{tab:all-selection-faithfulness}
    \begin{threeparttable}
    \begin{tabular}{llcc}
        \toprule
        Arm & Method & $\rho_\text{imp}$ ($q$) & $\rho_\text{dir}$ ($q$) \\
        \midrule
        Linear & SHAP/IG/GA     & $+0.377$ (0.0001)$^*$ & $-0.097$ (0.303) \\
        Linear & Probe          & $+0.289$ (0.003)$^*$  & $+0.399$ (0.0001)$^*$ \\
        \midrule
        MLP    & SHAP           & $+0.347$ (0.0004)$^*$ & $-0.091$ (0.473) \\
        MLP    & Probe saliency & $+0.122$ (0.343)      & n/a \\
        MLP    & IG             & $+0.437$ ($<$0.0001)$^*$ & $-0.058$ (0.630) \\
        MLP    & GA             & $+0.403$ ($<$0.0001)$^*$ & $-0.015$ (0.875) \\
        \bottomrule
    \end{tabular}
    \begin{tablenotes}\footnotesize
        \item[$^*$] Survives FDR 0.05. Linear arm: 5/8 tests survive (IG and GA counted separately from SHAP); MLP arm: 3/7. The same tests survive when BH is pooled over all 30 tests (both arms, both intervention scopes).
    \end{tablenotes}
    \end{threeparttable}
\end{table}

\clearpage
\section{Additional Limitations}
\label{app:limitations}

\begin{enumerate}
  \item \textbf{We calibrate $\alpha$ against a different readout than we
  evaluate it with} (Appendix~\ref{app:alpha}). We use a layer-11 residual
  probe's $P(\text{positive})$ at calibration time, and a GPT-2 logit
  difference at evaluation time.
  \item \textbf{The MLP arm's probe saliency score is magnitude-only}
  ($\|W_1[i,:]\|_2$, no sign). This is valid for importance/ranking comparisons,
  invalid for any directional or causal-sign comparison. We exclude it
  from group-steering vector construction for the same reason.
  \item \textbf{\texttt{top} selection biases SHAP's faithfulness $\rho$
  downward} relative to the other three methods, by range-restricting only
  SHAP's scored variance (Appendix~\ref{app:selection}). We never present
  a \texttt{top} number without this caveat.
  \item \textbf{The MLP probe's reported validation accuracy is mildly
  optimistic}: we select its inputs using a val-derived (label-free)
  frequency filter.
\end{enumerate}



